\documentclass[10pt,aspectratio=169]{beamer}
\usetheme{Madrid}
\usecolortheme{seahorse}

% ── 宏包 ──
\usepackage[utf8]{inputenc}
\usepackage{ctex}
\usepackage{amsmath,amssymb,amsthm}
\usepackage{tikz}
\usetikzlibrary{arrows.meta,positioning,calc}
\usepackage{booktabs}
\usepackage{graphicx}

% ── 定理环境 ──
\newtheorem{dfn}{定义}[section]
\newtheorem{thm}{定理}[section]
\newtheorem{prop}{命题}[section]
\newtheorem{rmk}{注}[section]

% ── 自定义 ──
\newcommand{\calL}{\mathcal{L}}
\newcommand{\calB}{\mathcal{B}}
\newcommand{\calP}{\mathcal{P}}
\newcommand{\calM}{\mathcal{M}}
\newcommand{\VV}{\mathbf{V}}
\DeclareMathOperator{\dom}{dom}
\DeclareMathOperator{\ran}{ran}

% ── 标题 ──
\title{从 Fréchet 导数到递归定义定理}
\subtitle{用 ZFC 严格构造高阶可导性}
\author{力学爱好者}
\date{\today}

\begin{document}

% ============================================================
\begin{frame}
  \titlepage
\end{frame}

% ── 目录 ──
\begin{frame}{目录}
  \tableofcontents
\end{frame}

% ============================================================
\section{背景：Fréchet 导数}
% ============================================================

\begin{frame}{赋范空间中的一阶可导}
  设 $X, Y$ 为赋范空间，$U \subseteq X$ 开集，$f: U \to Y$，$A \in U$。

  \medskip
  \begin{dfn}[Fréchet 可导]
    $f$ 在 $A$ 处 \textbf{可导}，若 $\exists\, L \in \calB(X, Y)$ 使得
    \[
      \lim_{\|h\|_X \to 0}
      \frac{\|f(A+h) - f(A) - L(h)\|_Y}{\|h\|_X} = 0
    \]
    记 $Df(A) := L$，称为 $f$ 在 $A$ 处的 \textbf{Fréchet 导数}。
  \end{dfn}

  \medskip
  \begin{itemize}
    \item $Df(A)$ 是 \textbf{有界线性算子}，不是数
    \item 映射 $Df: U \to \calB(X, Y)$，将每点映到其导数算子
    \item $\calB(X, Y)$ 本身是赋范空间（算子范数）
  \end{itemize}
\end{frame}

% ── 问题 ──
\begin{frame}{核心问题：如何定义 $n$ 阶导数？}
  自然的想法：反复求导。但严格地说——

  \medskip
  \begin{itemize}
    \item $f: U \to Y$，$Df: U \to \calB(X, Y)$
    \item $D^2f = D(Df): U \to \calB\bigl(X,\, \calB(X, Y)\bigr)$
    \item $D^3f: U \to \calB\Bigl(X,\, \calB\bigl(X,\, \calB(X, Y)\bigr)\Bigr)$
    \item \ldots 越来越深的嵌套！
  \end{itemize}

  \medskip
  \begin{block}{问题}
    \begin{enumerate}
      \item 这些嵌套算子空间 $\calL^n(X,Y)$ \textbf{存在且唯一}吗？
      \item $n$ 阶可导的定义如何\textbf{不依赖直觉}而严格表述？
      \item 这一切在 ZFC 集合论中是否\textbf{合法}？
    \end{enumerate}
  \end{block}
\end{frame}

% ============================================================
\section{ZFC 中的类与集合}
% ============================================================

\begin{frame}{ZFC 的基本设定}
  ZFC 的一阶语言：
  \begin{itemize}
    \item 只有一个二元谓词 $\in$
    \item 变量的论域 = \textbf{所有集合的全体}
    \item \textbf{论域中没有非集合的对象}
  \end{itemize}

  \medskip
  \begin{block}{``类''是什么？}
    ``类'' 不是 ZFC 的对象，而是\textbf{元语言简写}：
    \[
      \{x \mid \varphi(x)\} \quad \text{只是公式 } \varphi \text{ 的别名}
    \]
    \begin{itemize}
      \item 若 $\exists S\, \forall x\, (x \in S \leftrightarrow \varphi(x))$，则该类是\textbf{集合}
      \item 否则称为\textbf{真类} (proper class)，如 $\VV = \{x \mid x = x\}$
    \end{itemize}
  \end{block}
\end{frame}

\begin{frame}{``类函数''的含义}
  \begin{dfn}[类函数]
    一阶公式 $\psi(x, y)$ 若满足 $\forall x\, \exists! y\, \psi(x, y)$，
    则定义了一个\textbf{类函数} $G$：
    \[
      G(x) = y \;\iff\; \psi(x, y)
    \]
  \end{dfn}

  \begin{rmk}
    \begin{itemize}
      \item $G$ \textbf{不是} ZFC 论域中的对象——它只是公式 $\psi$
      \item 但对每个\textbf{集合}输入 $x$，输出 $G(x)$ 是唯一的集合
      \item 例：$G(Z) = \calB(X, Z)$（有界线性算子空间）
    \end{itemize}
  \end{rmk}
\end{frame}

\begin{frame}{NBG：类作为一等对象}
  在 NBG 集合论中，情况不同：

  \medskip
  \begin{columns}[T]
    \column{0.48\textwidth}
    \textbf{ZFC}
    \begin{itemize}
      \item 论域中只有集合
      \item ``$x$ 是集合''恒真
      \item 类 = 元语言缩写
      \item 类函数 = 一阶公式
    \end{itemize}

    \column{0.48\textwidth}
    \textbf{NBG}
    \begin{itemize}
      \item 论域包含类和集合
      \item $\mathrm{Set}(x) :\Leftrightarrow \exists y\, (x \in y)$
      \item 类是论域中的对象
      \item 类函数是类对象
    \end{itemize}
  \end{columns}

  \medskip
  \begin{block}{关键事实}
    NBG 是 ZFC 的\textbf{保守扩展}：关于集合的定理，两者证明力相同。
  \end{block}
\end{frame}

% ============================================================
\section{递归定义定理}
% ============================================================

\begin{frame}{Recursion Theorem on $\omega$}
  \begin{thm}[递归定义定理]
    设 $\psi(x,y)$ 是一阶公式满足 $\forall x\, \exists! y\, \psi(x,y)$（定义类函数 $G$）。
    给定任意集合 $a$，存在\textbf{唯一}的函数 $F: \omega \to \VV$（$F$ 是集合）使得
    \[
      F(0) = a, \qquad F(n+1) = G(F(n)) \quad \forall n \in \omega
    \]
  \end{thm}

  \begin{rmk}
    \begin{itemize}
      \item $G$ 可以是真类函数（只是公式），\textbf{不要求}是集合
      \item 但 $F$ \textbf{必须是集合}——因为 $\dom(F) = \omega$ 是集合
    \end{itemize}
  \end{rmk}
\end{frame}

\begin{frame}{证明思路：有限近似}
  对 $n \in \omega$，定义 $t_n: n+1 \to \VV$ 为 ``$n$\textbf{-近似}''：
  \[
    t_n(0) = a, \quad t_n(k+1) = G(t_n(k)) \;\; \forall k < n
  \]

  \medskip
  \textbf{证明分四步：}

  \medskip
  \begin{enumerate}
    \item \textbf{归纳}: 证 $t_n$ 存在且唯一（对 $n$ 归纳）
    \item \textbf{相容性}: $t_n \subseteq t_{n+1}$（嵌套链）
    \item \textbf{构造}: $F := \bigcup_{n \in \omega} t_n$（并集公理）
    \item \textbf{验证}: $F$ 满足递归条件且唯一
  \end{enumerate}
\end{frame}

\begin{frame}{证明细节——归纳步}
  \textbf{归纳假设}: $t_n$ 存在（是集合）。

  \medskip
  构造 $t_{n+1}$：
  \[
    t_{n+1} = t_n \cup \bigl\{\bigl(n+1,\; G(t_n(n))\bigr)\bigr\}
  \]

  \medskip
  为什么 $t_{n+1}$ 是集合？
  \begin{enumerate}
    \item $t_n(n)$ 是集合（从 $t_n$ 中取值，并集公理）
    \item $G(t_n(n))$ 是集合（$\forall x\, \exists! y\, \psi(x,y)$ 保证）
    \item 有序对 $(n+1, G(t_n(n)))$ 是集合（Kuratowski 定义 + 配对公理）
    \item 单元素集 $\{(n+1, G(t_n(n)))\}$ 是集合（配对公理）
    \item $t_n \cup \{\ldots\}$ 是集合（并集公理）
  \end{enumerate}
\end{frame}

\begin{frame}{证明细节——构造 $F$}
  \textbf{关键}: 序列 $(t_n)$ 是嵌套链：
  \[
    t_0 \subseteq t_1 \subseteq t_2 \subseteq \cdots
  \]

  \medskip
  由替换公理，$\{t_n \mid n \in \omega\}$ 是集合。

  定义
  \[
    F := \bigcup_{n \in \omega} t_n = \bigcup\, \{t_n \mid n \in \omega\}
  \]

  \begin{itemize}
    \item \textbf{$F$ 是集合}: 并集公理
    \item \textbf{$F$ 是函数}: 嵌套链的并保持函数性（若 $(k, y_1), (k, y_2) \in F$，
          则它们同属于某个 $t_N$，而 $t_N$ 是函数，故 $y_1 = y_2$）
    \item \textbf{$\dom(F) = \omega$}: 每个 $k$ 都在 $t_k$ 的定义域中
  \end{itemize}
\end{frame}

% ============================================================
\section{严格构造 $n$ 阶可导}
% ============================================================

\begin{frame}{构造步函数 $G$}
  定义\textbf{配置类}及其上的步函数：

  \medskip
  \begin{dfn}[配置]
    $(g, V, Z) \in \mathsf{Conf}_X$ 表示：$Z$ 是赋范空间，$V \subseteq X$ 开集，$g: V \to Z$。
  \end{dfn}

  \begin{dfn}[步函数 $G$]
    \[
      G(g, V, Z) = \begin{cases}
        (Dg,\; V,\; \calB(X, Z)) & \text{若 } g \text{ 在 } V \text{ 上处处 Fréchet 可导} \\
        \bot & \text{否则}
      \end{cases}
    \]
    其中 $\bot = \varnothing$ 表示 ``未定义''。
  \end{dfn}

  \begin{rmk}
    $G$ 由一阶公式定义，是合法的类函数。
  \end{rmk}
\end{frame}

\begin{frame}{应用递归定理}
  给定 $f: U \to Y$，取初始值 $a = (f, U, Y)$。

  \medskip
  由 Recursion Theorem，存在唯一的 $F: \omega \to \VV$：

  \[
    \begin{array}{rcl}
      F(0) &=& (f,\; U,\; Y) \\[4pt]
      F(1) &=& (Df,\; U,\; \calB(X,Y)) \\[4pt]
      F(2) &=& (D^2f,\; U,\; \calB(X, \calB(X,Y))) \\[4pt]
      &\vdots& \\
      F(n) &=& (D^n f,\; U,\; \calL^n(X,Y)) \quad \text{或 } \bot
    \end{array}
  \]

  \medskip
  $\bot$ 具有\textbf{吸收性}: 一旦 $F(k) = \bot$，则 $F(k+1) = G(\bot) = \bot$，此后全为 $\bot$。
\end{frame}

\begin{frame}{$n$ 阶可导的严格定义}
  \begin{dfn}[$n$ 阶可导（开集上）]
    $f$ 在 $U$ 上 $n$ 阶 Fréchet 可导 $\iff$
    \[
      \boxed{F_{f,U,Y}(n) \neq \bot}
    \]
    此时 $D^n f := \pi_1(F(n)): U \to \calL^n(X, Y)$。
  \end{dfn}

  \begin{dfn}[$n$ 阶可导（一点处）]
    $f$ 在 $A \in U$ 处 $n$ 阶可导 $\iff$
    \[
      \exists\, V \subseteq U \text{ 开},\; A \in V:\; F_{f|_V,\, V,\, Y}(n) \neq \bot
    \]
  \end{dfn}

  \begin{dfn}[光滑性]
    $f \in C^\infty(U, Y) \iff \forall n \in \omega:\; F(n) \neq \bot$
  \end{dfn}
\end{frame}

\begin{frame}{等价展开（非递归形式）}
  将递归展开，$F(n) \neq \bot$ 等价于：

  \medskip
  \begin{prop}
    $\exists$ 映射序列 $g_0, g_1, \ldots, g_n$ 和空间序列 $Z_0, Z_1, \ldots, Z_n$ 使得：
    \begin{enumerate}
      \item $g_0 = f$，$Z_0 = Y$
      \item 对每个 $k < n$：$Z_{k+1} = \calB(X, Z_k)$
      \item $g_k: U \to Z_k$ 在 $U$ 上处处可导
      \item $g_{k+1} = Dg_k$
    \end{enumerate}
  \end{prop}

  \medskip
  证明: 对 $n$ 归纳。\textbf{这正是递归定义定理的直接推论。}
\end{frame}

% ============================================================
\section{涉及的类为什么是集合}
% ============================================================

\begin{frame}{$\calL^n(X, Y)$ 是集合}
  \textbf{证明} (对 $n$ 归纳):

  \medskip
  \begin{itemize}
    \item[$n=0$:] $\calL^0 = Y$ 是集合（给定）
    \item[$n \to n+1$:] 设 $\calL^n$ 是集合。则
    \[
      \calL^{n+1} = \calB(X, \calL^n) \subseteq \calP(X \times \calL^n)
    \]
    \begin{itemize}
      \item $X \times \calL^n$ 是集合（配对 + 幂集）
      \item $\calP(X \times \calL^n)$ 是集合（幂集公理）
      \item $\calB(X, \calL^n)$ 是集合（分离公理：``$T$ 线性且有界''是一阶条件）
    \end{itemize}
  \end{itemize}

  \medskip
  进一步，整个序列 $(\calL^n)_{n \in \omega}$ 也是集合（Recursion Theorem 直接保证）。
\end{frame}

\begin{frame}{$C^\infty(U, Y)$ 是集合}
  \textbf{证明}:

  \medskip
  \begin{enumerate}
    \item 所有映射 $Y^U \subseteq \calP(U \times Y)$ 是集合
    \item ``$F_{f,U,Y}(n) \neq \bot$'' 可写为一阶公式 $\Theta(f, n)$
          （量词范围：有限序列 $g_0, \ldots, g_n$，均为集合）
    \item 因此
    \[
      C^\infty(U, Y) = \bigl\{ f \in Y^U \;\big|\; \forall n \in \omega\; \Theta(f, n) \bigr\}
    \]
    \item $\forall n \in \omega$ 量词范围是集合 $\omega$
    \item 由\textbf{分离公理}，$C^\infty(U, Y)$ 是集合 \qed
  \end{enumerate}
\end{frame}

\begin{frame}{配置类 $\mathsf{Conf}_X$ 是真类}
  $\mathsf{Conf}_X = \{(g, V, Z) \mid Z \text{ 是赋范空间}, V \subseteq X \text{ 开}, g: V \to Z\}$

  \medskip
  \textbf{这是真类}——因为 $Z$ 可取遍所有赋范空间，而 ``所有赋范空间'' 的收集是真类。

  \medskip
  \begin{block}{不影响证明}
    \begin{itemize}
      \item Recursion Theorem \textbf{不要求} $G$ 是集合
      \item 只要求 $G$ 是类函数（由一阶公式定义）
      \item 关键保证：$F$ 的每一步值 $F(n)$ 都是集合
    \end{itemize}
  \end{block}
\end{frame}

% ============================================================
\section{总结}
% ============================================================

\begin{frame}{逻辑链总结}
  \begin{center}
    \begin{tikzpicture}[
      scale=0.88, transform shape,
      node distance=1cm and 1.6cm,
      box/.style={draw, rounded corners, fill=blue!8, text width=2.6cm, align=center, font=\small},
      arr/.style={-Stealth, thick}
    ]
      \node[box] (zfc) {ZFC 公理系统 \\ (集合 + 一阶逻辑)};
      \node[box, right=of zfc] (rt) {Recursion Theorem \\ (有限近似 + 并集)};
      \node[box, right=of rt] (G) {步函数 $G$ \\ (求导操作)};
      \node[box, right=of G] (def) {$n$ 阶可导定义 \\ $F(n) \neq \bot$};

      \draw[arr] (zfc) -- node[above, font=\scriptsize]{提供公理} (rt);
      \draw[arr] (rt) -- node[above, font=\scriptsize]{保证迭代} (G);
      \draw[arr] (G) -- node[above, font=\scriptsize]{应用} (def);
    \end{tikzpicture}
  \end{center}

  \medskip
  \textbf{用到的 ZFC 公理:}
  \begin{columns}[T]
    \column{0.48\textwidth}
    \begin{itemize}
      \item 空集公理 ($0 = \varnothing$)
      \item 配对公理 (有序对)
      \item 并集公理 (构造 $F$)
    \end{itemize}
    \column{0.48\textwidth}
    \begin{itemize}
      \item 幂集公理 ($\calL^n$ 是集合)
      \item 分离公理 ($C^\infty$ 是集合)
      \item 替换公理 (序列存在)
    \end{itemize}
  \end{columns}
\end{frame}

\begin{frame}{两个框架的选择}
  \begin{columns}[T]
    \column{0.48\textwidth}
    \textbf{ZFC 视角}
    \begin{itemize}
      \item 类 = 公式的别名
      \item ``$x$ 是集合'' 恒真
      \item 真正要证的是\textbf{存在性}
      \item 更简洁，少一类对象
    \end{itemize}

    \column{0.48\textwidth}
    \textbf{NBG 视角}
    \begin{itemize}
      \item 类 = 论域中的对象
      \item ``$x$ 是集合'' 非平凡
      \item 表达更自然直观
      \item 保守扩展，不增加证明力
    \end{itemize}
  \end{columns}

  \medskip
  \begin{block}{核心 takeaway}
    无论选哪个框架，\textbf{递归定义定理}都能保证：
    从一阶导数出发，逐层构造出任意有限阶导数的严格定义——
    每一步都是 ZFC 公理合法的集合构造。
  \end{block}
\end{frame}

\begin{frame}{参考读物}
  \begin{thebibliography}{99}
    \bibitem{dieudonne} J. Dieudonné, \emph{Foundations of Modern Analysis}, Academic Press, 1969.
    \bibitem{cartan} H. Cartan, \emph{Differential Calculus}, Hermann, 1971.
    \bibitem{kunen} K. Kunen, \emph{Set Theory: An Introduction to Independence Proofs}, North-Holland, 1980.
    \bibitem{jech} T. Jech, \emph{Set Theory}, 3rd ed., Springer, 2003.
    \bibitem{mendelson} E. Mendelson, \emph{Introduction to Mathematical Logic}, 6th ed., CRC Press, 2015.
  \end{thebibliography}
\end{frame}

\end{document}
